\subsection{Algebraic extensions}

DEFINITION 1. --- An extension E of a field K is said to be algebraic (over K) if every element of E is algebraic over K. An extension E of K which is not algebraic is called transcendental (over K).

PROPOSITION 1. --- For an extension E of K to be algebraic it is necessary and sufficient that every sub-K-algebra A of E should be a field.

The condition is necessary: if E is algebraic over K and \(x \neq 0\) is an element of a sub-K-algebra A of E, then \(x^{-1} \in K[x] \subset A\) by V, p.~16, Th. 1, d). Therefore A is a field.

The condition is sufficient: if it is satisfied and x is an element \(\neq 0\) in E, then the ring K{[}x{]} is a field, hence \(x^{-1} \in K[x]\) ; in other words, there exists a polynomial \(g \in K[X]\) such that \(x^{-1} = g(x)\) , or also \(xg(x) - 1 = 0\) ; this shows x to be algebraic over K, hence E is an algebraic extension of K.

PROPOSITION 2. --- If an extension E of K is of finite degree n, then it is algebraic and the degree over K of each element of E divides n.

For, given \(x \in \mathbb{E}\) , \([\mathbf{K}(x): \mathbf{K}]\) is finite and divides \(n(\mathbf{V}, p. 10, \text{Cor. } 1)\) and hence \(x\) is algebraic over \(\mathbf{K}(\mathbf{V}, p. 17, \text{Cor. } 1)\) .

* There exist algebraic extensions of infinite degree, for example the algebraic closure of a finite field (V, p.~24, Remark 4). *

THEOREM 2. --- Let E be a finitely generated extension of K, generated by the elements \(a_1, \ldots, a_m\) which are algebraic over K; then E is an extension of finite degree of K. If the degree of \(a_i\) over K ( \(a_1, a_2, \ldots, a_{i-1}\) ) is \(n_i\) (for \(1 \leq i \leq m\) ), then the degree of E over K is \(n_1 n_2 \ldots n_m\) and the elements \(a_1^{\nu_1} a_2^{\nu_2} \ldots a_m^{\nu_m}\) ( \(0 \leq \nu_i \leq n_i - 1\) ) form a basis of E over K.

The elements \(a_i^{\nu_i}\) ( \(0 \leqslant \nu_i \leqslant n_i - 1\) ) form a basis of \(K(a_1, a_2, \ldots, a_i)\) over \(K(a_1, a_2, \ldots, a_{i-1})\) by Th. 1, b) of V, p.~16; the theorem therefore follows by induction on \(m\) from Prop. 25 of II, p.~222.

COROLLARY 1. --- Let E be an extension of K and A a subset of E consisting of elements that are algebraic over K. Then K(A) is algebraic over K and we have K{[}A{]} = K(A).

For each \(x \in \mathrm{K}(\mathrm{A})\) belongs to a field \(\mathrm{K}(\mathrm{F})\) , where F is a finite subset of A (V, p.~11, Cor.) ; now \(\mathrm{K}(\mathrm{F})\) is algebraic over K and equal to \(\mathrm{K}[\mathrm{F}]\) by Th. 2, hence x is algebraic over K and \(\mathrm{K}(\mathrm{A}) = \mathrm{K}[\mathrm{A}]\) .

COROLLARY 2. --- Let L be an extension of K and E, F subextensions of L. If F is algebraic over K, the subring K{[}E, F{]} of L generated by E ∪ F is a field coinciding with E(F) and algebraic over E.

For each element of F, being algebraic over K, is also algebraic over E (V, p.~17, Cor. 2) hence \(E(F)\) is an algebraic extension of E, and we have \(E(F) = E[F]\) by Cor. 1.

Remarks. --- 1) With the notation of Th. 2, E = K{[}a₁, a₂, \ldots, aₘ{]} and hence E is isomorphic to a quotient K{[}X₁, X₂, \ldots, Xₘ{]}/α; since E is a field, α is a maximal ideal in K{[}X₁, \ldots, Xₘ{]}.

\begin{enumerate}
\def\labelenumi{\arabic{enumi})}
\setcounter{enumi}{1}
\tightlist
\item
  Let \(\mathbf{E}\) be an algebraic extension of \(\mathbf{K}\) , of infinite degree. By Th. 2 there exists an infinite sequence \((a_n)_{n \geq 1}\) of elements of \(\mathbf{E}\) such that \(a_n \notin \mathbf{K}(a_1, a_2, \ldots, a_{n-1})\) ; Th. 2 shows further that the degree of \(\mathbf{K}(a_1, a_2, \ldots, a_n)\) over \(\mathbf{K}\) takes arbitrarily large values. In other words, if E is an algebraic extension of K such that the degrees \([F:K]\) of the subextensions F of E of finite degree over K are bounded, then E is an extension of finite degree of K.
\end{enumerate}

